\section{预算分档编排方法}
\label{sec:method}

本节描述四档协议在实现层面的组织方式。系统围绕同一个外层闭环运行：读入任务、生成候选、执行测试、记录结果；若失败，则在当前预算档位允许的范围内追加修补、重链、救援或级联。将程序生成与执行反馈连接起来的思想可追溯到执行引导程序综合和 REPL 式程序搜索\cite{chen2019execution,ellis2019write}。整体流程见图~\ref{fig:overall_framework}。

\begin{figure*}[t]
\centering
\includegraphics[width=0.98\textwidth]{gpt-fig2-overall-framework.png}
\caption{预算分档编排整体框架}
\label{fig:overall_framework}
\end{figure*}


\subsection{总体流程}
\label{sec:overview-flow}

对每个任务 $x_i$，系统首先生成单步候选并交由执行器判定。候选通过后流程结束；候选失败后，低档位直接结束或进入轻量修补，高档位根据结构难度、执行反馈与候选池状态决定后续动作。这样的组织把新增预算集中在\textbf{单步未能覆盖的残差样本}上。

四档的差异可以概括如下。单步档只做一次生成，是最低成本基线；单步+修补档在失败后追加有限修补调用，检验执行反馈的边际价值；分流+重链与救援档先判断样本是否进入重路径，再使用结构化提示、错误压缩与候选救援处理复杂失败；级联档进一步扩展候选集合，并用执行筛选和可选精修组织高预算搜索。

算法~\ref{alg:budget_tier} 给出单个样本上的预算分档编排逻辑。算法只刻画档位间的控制流，结构难度打分、失败反馈压缩与级联候选扩展的具体实现分别在后续小节展开。

\begin{algorithm*}[t]
\SetAlgoLined
\small
\caption{预算分档编排协议}
\label{alg:budget_tier}
\KwIn{任务 $x_i$，预算档位 $b$，模型 $p_\theta$，执行器 $E$，结构阈值 $\tau_{\mathrm{eff}}$}
\KwOut{最终候选 $y_i$，判定 $z_i$，调用次数 $c_i$，令牌消耗 $t_i$}
\KwWhere{$\textsc{Single}\prec\textsc{Repair}\prec\textsc{Route}\prec\textsc{Cascade}$ defines the budget order}
$y^\star \leftarrow \mathrm{Generate}(p_\theta, x_i)$\;
$(z^\star,e^\star) \leftarrow E(y^\star)$; update $c_i,t_i$\tcp*{single-step baseline}
\If{$z^\star=1$ \textbf{or} $b=\textsc{Single}$}{
    \Return $(y^\star,z^\star,c_i,t_i)$\;
}
\If{$b \succeq \textsc{Repair}$}{
    $\tilde{x}_i \leftarrow \mathrm{Concat}(x_i,\phi(e^\star))$\tcp*{compressed feedback}
    $y \leftarrow \mathrm{Generate}(p_\theta,\tilde{x}_i)$; $(z,e)\leftarrow E(y)$; update $c_i,t_i$\;
    \lIf{$z=1$}{$(y^\star,z^\star,e^\star)\leftarrow(y,z,e)$}
    \If{$z^\star=1$ \textbf{or} $b=\textsc{Repair}$}{
        \Return $(y^\star,z^\star,c_i,t_i)$\;
    }
}
$s_i \leftarrow \mathrm{Difficulty}(x_i)$; $r_i\leftarrow\mathbb{I}(s_i\ge\tau_{\mathrm{eff}})$\;
\If{$r_i=1$}{
    $(y,z,e)\leftarrow\mathrm{HeavyChain}(x_i,e^\star)$; update $c_i,t_i$\;
    \lIf{$z=1$}{$(y^\star,z^\star,e^\star)\leftarrow(y,z,e)$}
}
\If{$z^\star=0$}{
    $(y,z,e)\leftarrow\mathrm{Rescue}(x_i,e^\star)$; update $c_i,t_i$\;
    \lIf{$z=1$}{$(y^\star,z^\star,e^\star)\leftarrow(y,z,e)$}
}
\If{$b=\textsc{Cascade}$ \textbf{and} $z^\star=0$}{
    $\mathcal{Y}\leftarrow\mathrm{ExpandCandidates}(x_i,y^\star,e^\star)$\;
    $(y^\star,z^\star)\leftarrow\mathrm{SelectByExecutor}(\mathcal{Y},E)$; update $c_i,t_i$\;
}
\Return $(y^\star,z^\star,c_i,t_i)$\;
\end{algorithm*}

\subsection{结构难度路由}
\label{sec:routing}

难度路由用于在进入重路径之前完成一次粗粒度预算筛分。它承担\textbf{预算分配先验}的角色：结构上更复杂的题目更可能需要额外上下文、结构化提示或候选扩展，因而优先进入重路径。

图~\ref{fig:routing_decision} 展示了结构分数到轻/重路径的门控关系。

\begin{figure*}[t]
\centering
\includegraphics[width=0.98\textwidth]{gpt-fig3-routing-decision.png}
\caption{结构难度路由决策}
\label{fig:routing_decision}
\end{figure*}


记样本 $x_i$ 的结构难度分数为 $s_i$。实现中的 v2 打分器将分数拆为题面复杂度、签名/接口复杂度、测试复杂度、边界条件密度、题面歧义以及有界体量项：
\begin{equation}
\label{eq:difficulty}
s_i = p_i+\sigma_i+\tau_i+b_i+a_i+e_i.
\end{equation}
其中 $p_i,\sigma_i,\tau_i,b_i,a_i$ 分别对应题面、签名、测试、边界与歧义五类子项，$e_i$ 为由题面长度、测试长度、断言数量等构成的有界体量项。各子项由固定启发式规则加权并裁剪上限，总分由题面结构确定。

路由决策写为
\begin{equation}
    r_i=\mathbb{I}(s_i\ge \tau_{\mathrm{eff}}),
\end{equation}
其中 $\tau_{\mathrm{eff}}$ 为配置中的结构阈值，主实验第三档默认使用 $1.6$。$r_i=1$ 表示进入重路径主链，$r_i=0$ 表示停留在轻路径。阈值越低，重路径覆盖率越高，通常带来更高成本和更高通过率上限；阈值越高，成本下降，同时覆盖样本减少。该关系由第~\ref{sec:experiments} 节的阈值消融检验。

为分离重路径覆盖率与结构对齐的影响，仓库实现了 \texttt{random\_permute\_matched} 对照：先按结构分得到本批样本的轻/重路径标签多重集，再固定随机种子置换标签，使重路径数量与默认结构路由完全一致，但题目与重路径槽位的对应被打乱。该对照刻画固定覆盖率下的结构对齐效果。

\subsection{失败反馈与救援链}
\label{sec:rescue}

执行失败后，系统把失败信息转化为下一轮生成可利用的局部约束。其基本思路与 Self-Debug、Reflexion 等执行反馈工作一致\cite{chen2024selfdebug,shinn2023reflexion}，也与利用工具反馈进行自我修正的研究方向相近\cite{gou2024critic,madaan2023selfrefine}。本文进一步考察反馈形式对令牌成本与修补收益的影响。

在分流+重链与救援档中，系统先根据执行器返回的 \texttt{check\_correctness} 字符串判断失败形态。该判断由规则函数完成：未闭合字符串、语法非法等更倾向于替换生成；缺参、参数数量不匹配、名称未定义或缩进错误等归入修补；其余进入更保守的替换或救援路径。当单步与结构化双路径均失败时，系统比较二者错误标签，以决定后续修补轮数。

压缩后的回注包记为
\begin{equation}
\begin{aligned}
\phi(e_i^{(1)})=\operatorname{Concat}\bigl(&
\textit{错误类型},
\textit{截断摘要},\\
&
\textit{规则片段},
\textit{定向修补提示}
\bigr).
\end{aligned}
\end{equation}
其中 $e_i^{(1)}$ 为第一次失败反馈。第二轮上下文可写作
\begin{equation}
\tilde{x}_i^{(2)}=\operatorname{Concat}(x_i,\phi(e_i^{(1)})),
\end{equation}
并据此生成新候选
\begin{equation}
y_i^{(2)} \sim p_{\theta}(\cdot\mid \tilde{x}_i^{(2)}).
\end{equation}
该表达式说明失败信息如何被裁剪并进入下一轮上下文。图~\ref{fig:rescue_loop} 展示了这一闭环。

\begin{figure}[t]
\centering
\includegraphics[width=0.98\linewidth]{gpt-fig4-rescue-loop.png}
\caption{失败反馈救援闭环}
\label{fig:rescue_loop}
\end{figure}


本文在实验中比较压缩回注与完整回溯两种形态。五随机种子结果显示，压缩回注保持相近通过率，同时平均令牌更低，因此作为默认回注形式。

\subsection{高预算级联}
\label{sec:cascade}

级联档把新增预算转向候选空间覆盖。其流程包括三层：候选扩展层生成额外候选，执行筛选层用统一测试选择可行实现，后处理层对进入候选池的样本做有限精修或融合。该设计与测试时扩展的基本思想一致，即通过候选扩展提高命中可行解的概率\cite{li2022alphacode,li2025sstar}。

实现上，级联主链按基线候选、额外候选扩展、执行筛选和可选后处理的顺序运行。额外候选阶段的生成次数受全局上限、路由表与执行 top-$k$ 共同约束；实际执行数为生成数与 top-$k$ 的较小值。可行解指在统一执行器下通过官方单元测试的候选实现。若候选扩展已产生可行解，精修只作为精度--成本微调。

第~\ref{sec:experiments} 节的级联消融显示，额外候选扩展是高预算档的主增益来源；精修相关开关更多体现为令牌与少量精度之间的折中。级联档的核心是用候选覆盖换取更高上限。

\subsection{实现不变式与可比性控制}
\label{sec:method-invariants}

为保证不同预算档位之间的比较只反映编排动作差异，本文在实现中固定三类不变式。第一，所有档位共享同一执行器、测试接口、超时设置和通过判定函数，避免把评测口径差异混入预算收益。第二，结构难度分数只由题面、签名和测试结构等预执行信息计算，不使用执行结果或人工标签；随机置换路由对照进一步把重路径覆盖率与结构对齐效果分开。第三，失败回注只使用执行器返回的错误信息及规则化摘要，不引入参考答案或额外人工修正。上述约束使四档协议能够作为同一测试时计算空间中的可记账比较，而不是彼此独立的提示模板集合。
